% !TeX encoding = UTF-8
% !TeX spellcheck = de_DE
% Datei: 03_variationsrechnung.tex
% Stabilitätsanalyse des Systems S mittels Variationsrechnung

\chapter{Die Soziale Variationsrechnung -- Stabilitätsanalyse}
\label{chap:variation}

\section{Einleitung: Methodentransfer von der Physik}
Die bisherige Darstellung des Systems $\mathcal{S}$ war normativ-logischer Natur. Um jedoch seine praktische Lebensfähigkeit zu beurteilen, muss seine Stabilität gegenüber Störungen analysiert werden. Hierzu wird das mathematische Instrumentarium der \emph{Variationsrechnung} (Calculus of Variations) adaptiert.

In der Physik sucht die Variationsrechnung nach Funktionen, die ein gegebenes Funktional (z.B. die Wirkung) extremalisieren, und untersucht die Stabilität dieser Lösungen gegen kleine Variationen. Übertragen auf die $\mathcal{S}$-Gesellschaft:
\begin{itemize}
    \item Das zu optimierende \emph{Funktional} ist der Demokratiegrad $\mathcal{D}(G)$.
    \item Die \emph{Variablen} sind gesellschaftliche Parameter wie die Reichtumsgrenze $\varepsilon$ oder die Bevölkerungsgröße $N$.
    \item Ein \emph{stabiles Gleichgewicht} liegt vor, wenn kleine Störungen dieser Parameter nur zu kleinen Änderungen in $\mathcal{D}(G)$ führen und Rückstellkräfte existieren.
\end{itemize}
Dieses Kapitel formalisiert diese Analogie und leitet zentrale Stabilitätsbedingungen her.

\section{Formale Grundlagen: Das Demokratizitäts-Funktional}

\begin{definition}[Demokratizitäts-Funktional]
    Für eine Gesellschaft $G$, die das System $\mathcal{S}$ instanziiert, definieren wir das \emph{Demokratizitäts-Funktional} $\mathfrak{D}$ als Abbildung der gesellschaftlichen Parameter auf den Demokratiegrad:
    \[
    \mathfrak{D}: \Theta \rightarrow [0,1], \quad \theta \mapsto \mathcal{D}(G_\theta)
    \]
    wobei $\Theta$ der Raum aller zulässigen Parameterkombinationen unter der Verfassung $C$ ist und $G_\theta$ die konkrete Ausprägung der Gesellschaft mit Parametern $\theta$ bezeichnet.
\end{definition}

\section{Modul 1: Analyse des Parameters $\varepsilon$ (Reichtumsgrenze)}
\label{sec:modul_epsilon}

\subsection{Modellannahmen und Randbedingungen}
\begin{enumerate}
    \item \textbf{Initialverteilung:} Die natürliche Verteilung von Produktivität, Glück und Präferenzen sei durch die Dichtefunktion $f(x)$ über ökonomischem Potenzial $x \geq 0$ beschrieben.
    \item \textbf{Umverteilungsmechanismus:} Die in Axiom~\ref{ax:3} geforderte Begrenzung wird durch eine 100\%-Steuer auf alle Anteile $x > \varepsilon$ modelliert, gefolgt von einer gleichmäßigen Redistribution. Die resultierende \emph{beschnittene Verteilung} ist $f_\varepsilon(x)$.
    \item \textbf{Konversionsannahme:} Ökonomisches Potenzial $x$ kann mit einer Rate $c > 0$ in politischen Einfluss $\iota(x)$ umgewandelt werden: $\iota(x) = c \cdot x$ für $x \leq \varepsilon$, und $\iota(x) = 0$ für $x > \varepsilon$ (da überschüssiges Vermögen abgeschöpft wird).
\end{enumerate}

\subsection{Variation und Instabilitätsszenarien}

\begin{theorem}[Instabilität bei zu hohem $\varepsilon$]
    Für $\varepsilon \to \infty$ (de-facto Abschaffung der Grenze) gilt:
    \[
    \lim_{\varepsilon \to \infty} \mathfrak{D}(\varepsilon) = 0.
    \]
    Das System kollabiert zu einer Oligarchie.
\end{theorem}

\begin{proof}
    Mit $\varepsilon \to \infty$ wird $f_\varepsilon(x) = f(x)$. Die Varianz $\operatorname{Var}(f)$ ist maximal. Die politische Einflussverteilung $\iota(x)$ hat daher ebenfalls maximale Varianz, was direkt $\mathcal{D}(G) \to 0$ impliziert (da politische Gleichheit nicht mehr gegeben ist). Zudem kann sich eine dauerhafte Elite etablieren.
\end{proof}

\begin{theorem}[Instabilität bei zu niedrigem $\varepsilon$]
    Sei $\varepsilon_{\text{sub}}$ ein Niveau nahe dem Subsistenzminimum. Dann gilt:
    \[
    \lim_{\varepsilon \to \varepsilon_{\text{sub}}} \mathfrak{D}(\varepsilon) = \delta \ll 1.
    \]
    Das System kollabiert in Stagnation oder informelle Kleptokratie.
\end{theorem}

\begin{proof}
    Bei sehr niedrigem $\varepsilon$:
    \begin{enumerate}
        \item Anreize für produktive und innovative Tätigkeiten verschwinden $\Rightarrow$ gesamtgesellschaftliche Produktion $Y(G)$ sinkt.
        \item Der für Bildung und demokratische Infrastruktur notwendige Überschuss schwindet, was $B(G) < B_{\min}$ zur Folge hat (vgl. Theorem 2 aus Kapitel~\ref{chap:theoreme}).
        \item Knappheit begünstigt die Entstehung \emph{informeller} Machtstrukturen um die Verteilung der Grundgüter (Nepotismus, Korruption), die undemokratisch und intransparent sind.
    \end{enumerate}
    Beide Effekte reduzieren $\mathfrak{D}$ drastisch.
\end{proof}

\subsection{Das Theorem des stabilen Intervalls}

\begin{theorem}[Existenz eines stabilen Intervalls für $\varepsilon$]
    Unter den gegebenen Annahmen existiert ein endliches, positives Intervall $I_{\text{eq}} = [\varepsilon^* - \delta, \varepsilon^* + \delta]$ mit $\varepsilon^* > \varepsilon_{\text{sub}}$, so dass:
    \[
    \boxed{
        \forall \varepsilon \in I_{\text{eq}}: \quad \mathfrak{D}(\varepsilon) \text{ ist maximal und stabil, d.h. } \frac{d\mathfrak{D}}{d\varepsilon} = 0 \text{ und } \frac{d^2\mathfrak{D}}{d\varepsilon^2} < 0.
    }
    \]
    Außerhalb von $I_{\text{eq}}$ wird $\mathfrak{D}$ monoton abfallen und das System instabil.
\end{theorem}

\begin{proof}[Beweisskizze]
    Die Funktion $\mathfrak{D}(\varepsilon)$ ist:
    \begin{enumerate}
        \item stetig in $\varepsilon$ (da Umverteilung stetig wirkt),
        \item $\mathfrak{D}(\varepsilon_{\text{sub}}) = \delta$ (nahe 0),
        \item $\mathfrak{D}(\varepsilon \to \infty) \to 0$,
        \item positiv für mittlere $\varepsilon$.
    \end{enumerate}
    Nach dem Zwischenwertsatz und den Extremwertbedingungen stetiger Funktionen auf kompakten Intervallen existiert daher mindestens ein lokales Maximum $\varepsilon^*$. Die zweite Ableitung muss an diesem Punkt negativ sein, sonst läge kein Maximum vor. Die Ausdehnung $\delta$ definiert den Bereich, in dem $\mathfrak{D}$ innerhalb einer tolerierten Abweichung vom Maximum bleibt.
\end{proof}

\subsection{Praktische Implikationen}
Dieses Theorem hat eine tiefgreifende Konsequenz: Die \emph{genaue wissenschaftliche Bestimmung von $\varepsilon^*$ und die präzise Einhaltung dieser Grenze} werden zur Überlebensfrage der $\mathcal{S}$-Gesellschaft. Dies stellt hohe Anforderungen an die wissenschaftliche Infrastruktur (Axiom~\ref{ax:4}) und wirft die Frage nach dem konkreten institutionellen Mechanismus zur Festsetzung von $\varepsilon$ auf.

\section{Modul 2: Analyse des Parameters $N$ (Bevölkerungsgröße)}
\label{sec:modul_n}

\subsection{Das Skalierungsproblem direkter Demokratie}

\begin{theorem}[Kritische Größe für direkte Deliberation]
    Sei $t_{\text{delib}}$ die durchschnittliche Zeit, die für die informierte Deliberation einer einzelnen politischen Frage pro Person benötigt wird. Sei $n_q$ die Anzahl der zu entscheidenden Fragen pro Zeiteinheit. Dann existiert eine \emph{kritische Bevölkerungsgröße} $N_c$, gegeben durch:
    \[
    N_c = \frac{T_{\text{verfügbar}}}{t_{\text{delib}} \cdot n_q},
    \]
    wobei $T_{\text{verfügbar}}$ die pro Person für politische Partizipation verfügbare Gesamtzeit ist. Für $N > N_c$ ist reine direkte Demokratie unmöglich.
\end{theorem}

\subsection{Stabilität durch Netzwerk-Delegation}

\begin{theorem}[Stabile Delegationsnetzwerke]
    Für $N > N_c$ muss das System $\mathcal{S}$ in ein mehrstufiges, aber nicht-hierarchisches Delegationsnetzwerk übergehen, um stabil zu bleiben. Die Stabilität eines solchen Netzwerkes hängt von seiner Topologie ab. Sei $d_{\text{max}}$ der maximale Abstand (in Delegationsstufen) zwischen einem beliebigen Individuum und der finalen Entscheidungsebene. Dann gilt die \emph{Distanz-Stabilitäts-Bedingung}:
    \[
    \boxed{d_{\text{max}} \leq d_{\text{krit}} \quad \text{mit} \quad d_{\text{krit}} \approx -\frac{\ln(1 - \alpha)}{\ln(k)} }
    \]
    wobei $k$ der durchschnittliche Verzweigungsgrad (Delegationsverhältnis) und $\alpha$ die akzeptable Verzerrung des delegierten Willens ist. Wird $d_{\text{max}} > d_{\text{krit}}$, kollabiert das System in eine de-facto Hierarchie.
\end{theorem}

\section{Modul 3: Analyse des Parameters $\sigma$ (Informationsasymmetrie)}
\label{sec:modul_sigma}

\begin{definition}[Informationsasymmetrie-Maß]
    Sei $I_i$ das Informationsniveau des Individuums $i$ in relevanten politischen Sachfragen. Dann ist:
    \[
    \sigma(G) = \sqrt{\frac{1}{|I|} \sum_{i \in I} (I_i - \bar{I})^2}, \quad \bar{I} = \frac{1}{|I|} \sum_{i \in I} I_i.
    \]
\end{definition}

\begin{theorem}[Bildung als Stabilisator]
    Die zeitliche Entwicklung von $\sigma(G)$ gehorcht näherungsweise der Differentialgleichung:
    \[
    \frac{d\sigma}{dt} = \underbrace{\gamma_{\text{nat}}}_{\text{nat. Diffusionsrate}} - \underbrace{\beta \cdot B(G)}_{\text{Bildungseffekt}}.
    \]
    Dabei ist $\gamma_{\text{nat}} > 0$ eine Konstante, die die natürliche Entstehung von Wissensunterschieden beschreibt, und $\beta > 0$ die Effizienz des Bildungssystems. Für Stabilität ($\frac{d\sigma}{dt} \leq 0$) muss daher gelten:
    \[
    \boxed{B(G) \geq B_{\min} = \frac{\gamma_{\text{nat}}}{\beta}.}
    \]
    Dies bestätigt und quantifiziert Theorem 2 aus Kapitel~\ref{chap:theoreme}.
\end{theorem}

\section{Synthese: Vulnerabilitäten und Vorhersagen des Systems $\mathcal{S}$}

Die Variationsrechnung offenbart drei kritische Vulnerabilitäten der $\mathcal{S}$-Gesellschaft:
\begin{enumerate}
    \item \textbf{Die $\varepsilon$-Sensitivität:} Das System ist nur in einem schmalen Korridor ökonomischer Gleichheit stabil. Die Bestimmung und Durchsetzung der exakten Grenze $\varepsilon^*$ erfordert eine präzise, unbestechliche und hochautoritative wissenschaftliche Instanz -- eine potenzielle Keimzelle für Technokratie.
    \item \textbf{Das Skalierungsdilemma:} Jenseits der kritischen Größe $N_c$ muss das System komplexe Delegationsnetzwerke einführen, deren Stabilität von schwer zu kontrollierenden topologischen Parametern ($d_{\text{max}}$, $k$) abhängt.
    \item \textbf{Die Bildungsabhängigkeit:} Das System kollabiert deterministisch, wenn die Bildungsinvestition $B(G)$ unter die kritische Schwelle $B_{\min}$ fällt. Dies macht es anfällig für kurzfristige Krisen, die zu Budgetkürzungen führen könnten.
\end{enumerate}

\emph{Vorhersage 1:} Eine $\mathcal{S}$-Gesellschaft wird tendenziell klein oder stark dezentralisiert sein ($N \approx N_c$), um das Skalierungsdilemma zu umgehen.

\emph{Vorhersage 2:} Der Haushalt für Bildung und Wissenschaft wird in einer $\mathcal{S}$-Gesellschaft der mit Abstand größte und am stärksten geschützte Posten sein.

\emph{Vorhersage 3:} Der demokratischste Moment könnte gleichzeitig der technokratischste sein: die jährliche, evidenzbasierte Neufestsetzung der Reichtumsgrenze $\varepsilon$.

\section{Schlussfolgerung}
Die Variationsrechnung transformiert die normative Theorie $\mathcal{S}$ in ein \emph{dynamisches Modell} mit testbaren Stabilitätsbedingungen. Sie zeigt, dass die $\mathcal{S}$-Gesellschaft theoretisch stabil existieren kann, aber nur unter sehr spezifischen und streng kontrollierten Bedingungen. Ihre praktische Realisierung wäre ein permanentes Balancieren auf einem schmalen Grat zwischen den Abgründen der Oligarchie, der Stagnation und der Technokratie. Damit liefert diese Analyse den Übergang zur Diskussion der praktischen und philosophischen Herausforderungen im nächsten Kapitel.
